% ============================================================================
%  Chapitre 6, partie I : Limites du cas quasi idéal
% ============================================================================

% ----------------------------------------------------------------------------
\section{Limites du cas quasi idéal}
% ----------------------------------------------------------------------------

Le chapitre précédent repose sur des hypothèses qui ne tiennent pas sur
une vraie boule :
\begin{itemize}
    \item \textbf{le biais du gyromètre est constant} : en réalité, il
    dérive. $\vect{\gamma}$, intégré avec le gyromètre seul, dérive alors
    sans que rien ne puisse le corriger, et les régresseurs sont faussés.
    \item \textbf{l'accéléromètre ne sert qu'au repos} : pour recaler
    $\vect{\gamma}$ pendant le mouvement, il faut savoir ce qu'il mesure
    quand la boule roule, et ce n'est plus seulement $g\,\vect{\gamma}$.
\end{itemize}
On garde en revanche l'hypothèse de Chaplygin $G = C$ : sur une boule réelle, elle n'est vérifiée qu'approximativement, et l'effet d'un décalage de $G$ n'est pas étudié ici.
